\phantomsection\bheading{section}{6.\quad Errors}{section}{\protect\numberline{6.}Errors}

\phantomsection\bheading{subsection}{6.1\quad What a message must do}{subsection}{\protect\numberline{6.1}What a message must do}

Three places, a consequence, a fix, and a verdict on whether this is sharing or a limit. The three places are Rust's three points — the borrow, the invalidating action, the later use — with the blame on the action: “the error, conceptually, is to perform an invalidating action in between two uses of the reference” (RFC 2094, \emph{Leveraging intuition}). Futhark names the two sites and offers \code{copy} on every entry of its error index (§8.1.1–§8.1.13), and its 2021 lesson is that unnamed intermediate results must be given a description the user can read: “Consuming result of applying “iota” (at 2:23-28)” (2021 post, l. 236). Hylo draws the arrow from the consuming site to the later use and “will offer to insert these copies for you” (tour, l. 60). Czaplicki's two rules apply unchanged: the code as written, above the message; a hint under it, every time (2015). Wrenn and Krishnamurthi's frame — a message is a classifier of code into “look here” and “do not look here”, and its highlights should be exactly the places an edit could go (2017 §2, §8.4) — fixes the highlight set: the kept place, the edit, the later use, nothing else. Marceau's warning that highlights have “an over-focusing effect” and that “expected … found” wording “suggests that the definition is fine and the use problematic” (2011 §7.2) is why the message blames none of the three alone.

Crichton's finding is the one that shapes the verdict line: learners “could usually predict why the borrow checker would reject a program” (64–78 \%) “but could only fix the program in 46 \% of cases … and could only create a counterexample in 31 \%”, and “only 3/15 participants” saw that a rejected program was in fact safe (2023 §2.4). A sound-and-incomplete checker's errors “may arise from either ownership-unsound behavior or limitations of the analyzer. Understanding this distinction is essential for fixing ownership errors” (2020 abstract). So every message says which it is, and how it knows.

\phantomsection\bheading{subsection}{6.2\quad Examples}{subsection}{\protect\numberline{6.2}Examples}

\textbf{1. A later read (real sharing).}

\begin{Verbatim}
-- EDIT OF A LIST STILL IN USE ------------------------------------ src/Main.beni

`List.push` changes a list in place, but this one is still needed afterwards.

12|     saved = model.todos
                ^^^^^^^^^^^ `saved` keeps the list here
15|     next = List.push model.todos todo
                ^^^^^^^^^^^^^^^^^^^^^^^^^^ the list is changed in place here
19|     List.length saved
                    ^^^^^ and `saved` is read here, after the change

If the change happened in place, `saved` would already contain `todo` when it
is read on line 19: `List.length saved` would be one more than before.

To keep `saved` as it was, copy it:

    saved = List.copy model.todos

Or, if `saved` only needs the list before the change, use it before line 15.

This is sharing the checker can see: the list is read on line 19 on every path
through this function.
\end{Verbatim}

\textbf{2. A closure that keeps the list (real sharing).}

\begin{Verbatim}
`List.set` changes a list in place, but a function still holds it.

 8|     show = λi → String.fromInt (List.length rows)
                                                ^^^^ `show` holds `rows` here
11|     rows2 = List.set rows 0 r
                ^^^^^^^^^^^^^^^^^ the list is changed in place here
14|     { model | label = show 3 }
                          ^^^^ `show` is called here, after the change

`show` would see the changed list. If that is intended, call `show` with a
copy: `show = λi → String.fromInt (List.length (List.copy rows))` — or compute
what `show` needs before line 11.
\end{Verbatim}

\textbf{3. A command that captured the list (real sharing, across dispatches).}

\begin{Verbatim}
`List.push` changes `model.todos` in place, but a command started by another
message may still be reading it.

-- in the arm for `Save`:
41|     Cmd.perform (λsend → send (Saved (encode model.todos)))
                                                 ^^^^^^^^^^^ the command keeps `model.todos`
-- in the arm for `Add`:
27|     { model | todos = List.push model.todos todo }
                          ^^^^^^^^^^^^^^^^^^^^^^^^^^ changed in place here

A command runs after `update` returns and may run after later messages, so the
list it keeps is the same one `Add` changes. To give the command its own copy:

41|     Cmd.perform (λsend → send (Saved (encode (List.copy model.todos))))

Or encode before the command is made:

40|     text = encode model.todos
41|     Cmd.perform (λsend → send (Saved text))
\end{Verbatim}

\textbf{4. A once-closure called many times (real sharing).}

\begin{Verbatim}
This function changes `acc` in place, but `List.map` may call it once per element.

 7|     List.map xs (λx → List.push acc x)
                     ^^^^^^^^^^^^^^^^^^^^ changes `acc` in place
        ^^^^^^^^ `List.map` calls it for every element of `xs`

Each call would change the same list, so the result of the first call would
change under the second. To build a list by adding to `acc`, fold:

 7|     List.foldl xs acc (λx acc → List.push acc x)
\end{Verbatim}

\textbf{5. A checker limit.}

\begin{Verbatim}
`List.set` changes `rows` in place, but the checker cannot tell whether anything
else still holds it.

14|     rows = Js.to raw
               ^^^^^^^^^ `rows` came from JavaScript here, which may keep it
17|     List.set rows i v
        ^^^^^^^^^^^^^^^^^ changed in place here

This is a limit of the checker, not proof of sharing: nothing in this module
reads the old `rows` after line 17, but the checker cannot see what JavaScript
does with the array it gave you. To be safe, copy it once where it arrives:

14|     rows = List.copy (Js.to raw)
\end{Verbatim}

\phantomsection\bheading{subsection}{6.3\quad The shape, fixed}{subsection}{\protect\numberline{6.3}The shape, fixed}

\begin{Verbatim}
<TITLE> ------------------------------------------------------------ <file>

<one sentence: what primitive, what list, why that is a problem here>

<the kept place, in source order, with a label beginning "… keeps"/"… holds">
<the edit, with the label "changed in place here">
<the later use, with the label "… read/called here, after the change">

<one or two sentences: what the program would see if the change happened in place>

<the fix, machine-applicable, as a replacement line, "To … , copy it: …">
<an alternative fix that avoids the copy, when there is one>

<the verdict line: "This is sharing the checker can see: …" or "This is a limit
 of the checker, not proof of sharing: …", naming the limit>
\end{Verbatim}

The suggestion is \code{MachineApplicable} in rustc's vocabulary: the compiler can insert \code{List.copy} at the kept place by itself (rustc-dev-guide, \emph{Suggestions}). “Did you mean” is not used; the word “illegal” is not used.

\phantomsection\bheading{subsection}{6.4\quad Telling a checker limit from real sharing: the derivation tag}{subsection}{\protect\numberline{6.4}Telling a checker limit from real sharing: the derivation tag}

Every fact the checker derives carries a tag: \textbf{exact} when it comes from a read the source shows (a variable occurrence, a field access, a call whose summary was inferred from a beni body, a capture in a lambda written here), or \textbf{approximate} when it comes from one of the checker's own limits:

{\small\setlength{\emergencystretch}{3em}
\begin{longtable}{>{\raggedright\arraybackslash}p{\dimexpr0.350\linewidth-2\tabcolsep\relax}>{\raggedright\arraybackslash}p{\dimexpr0.650\linewidth-2\tabcolsep\relax}}
\toprule
\textbf{Limit} & \textbf{Tag text} \\
\midrule
\endfirsthead
\toprule
\textbf{Limit} & \textbf{Tag text} \\
\midrule
\endhead
\bottomrule
\endfoot
a \code{foreign} or \code{Js} with a declared $\mathsf{shared}$ or no summary & “came from / was given to JavaScript, which may keep it” \\[0.35em]
a cell marked $\top$ by the k-cut or a width cap & “the value is nested deeper than the checker follows (8 steps)” \\[0.35em]
a path-insensitive join (§5.6's second shape) & “the list may be either $a$ or $b$ here; the checker treats it as both” \\[0.35em]
a \code{Ref}/\code{Queue} read & “was taken from a \code{Ref}, whose other readers the checker cannot see” \\[0.35em]
a method whose summary is pessimistic (an unannotated platform \code{foreign}) & “\code{M.f} declares nothing about what it keeps” \\[0.35em]
$\mathit{calls}\ \mathsf{many}$ for a closure the program calls once in fact & “\code{List.map} may call it more than once” \\[0.35em]
\bottomrule
\end{longtable}}

A rejection whose derivation uses only exact tags is \textbf{real sharing}: there is a path through the program text on which the old version is read after the edit. One that uses an approximate tag is \textbf{possible sharing}, and the message names the limit. This is Crichton's distinction made mechanical, and it is also what §9's experiment counts. A third thing the tags give for free is Swift's “compiler bug” sentinel (DiagnosticsSIL.def, ll. 987–996): a rejection with no derivation at all is an internal error, never a user error.

\phantomsection\bheading{subsection}{6.5\quad The stability rule}{subsection}{\protect\numberline{6.5}The stability rule}

\begin{quote}
A program the checker of release \emph{n} accepts is accepted by every release \emph{m > n}.
\end{quote}

The accepted set is defined by the rules of §2–§3, not by what any optimiser removes (McCall: “these guarantees would have to be real guarantees … even in debug builds”; Farvardin: checking must happen “earlier in the optimization pipeline, just to guarantee stability across different versions”). Precision improvements are additive — a finer path language, a container function that gains an exact summary, a once-closure that gains a $\mathit{calls}\le 1$ position — and each only removes approximate tags. The one legitimate exception is a \textbf{soundness fix}: a rule found unsound (Futhark's 2026 self-aliased abstract type is the model) may reject previously accepted programs, and must do so with a diagnostic that says so and a \code{beni fmt --migrate-…} where a mechanical rewrite exists. The rule set is versioned in the interface format (a summary block version), so a cached interface from an older rule set is a miss, never a silent mix.
